\begin{textsample}{NEG Dim 2 – global\_north\_2024\_12 – Score 1.00 – global\_north\_2024\_12/119484401.txt}  \label{ex:f2_neg_167}
Hossam Madhoun worked for 30 years as a theatre-maker and writer in Gaza City before Israel began its devastating onslaught on the Palestinian enclave in October 2023."The theatre movement in Gaza is deprived of exchanging with the world, especially in the last 17 years with the Israeli blockade,"he tells Middle East Eye from Cairo, where he now lives with his family.
Since then, it has been more challenging, but the show went on : in 2009, he began collaborating with Az Theatre, a company based in London and run by Jonathan Chadwick, launching workshops and productions together.
Even then, Madhoun could not sustain himself purely as a playwright and many of his theatre projects now lie neglected as the war rages on.
New MEE newsletter : Jerusalem Dispatch Sign up to get the latest insights and analysis on Israel-Palestine, alongside Turkey Unpacked and other MEE newsletters Since 2013, he has worked as a child protection manager for Ma'an Development Centre, a Gaza-based NGO that provides psychosocial support, counselling, recreational activities and case management for child protection concerns."The most terrible thing children are going through is fear,"Madhoun says."Who can bear that?"' I don't know why Palestinians are not considered human or have to prove that they are human in order to be protected ' Mainstream media reports estimate that more than 45,500 people have been killed, including 17,000 children, but the count has been frozen at this number for months.
More than 10,000 are missing and presumed dead.
Born in the heat of war, dead from the cold : Gaza ' s children are freezing to death The real death toll could be more than quadruple that figure.
Estimates suggest that over 186,000 people are dying indirectly from hunger, cold, disease as well as the bombs that rain over them continually.
Like Madhoun ' s mother.
She died earlier this year from kidney failure, which could have been prevented, except there was no access to dialysis because there were no hospitals, because there were far younger, far more urgent cases that had to be prioritised with limited resources.
These people are not counted as victims of the war, but why else did they die?"I don't know why Palestinians are not considered human or have to prove that they are human in order to be protected,"Madhoun asks. ' A place between Earth and the sky ' In 2022, Madhoun collaborated with Az Theatre on an adaptation of Slawomir Mrozek ' s play The Emigrants, which is about those who left Gaza and the other side of immigration : how it can be devastating for people who have given up their land, their homes and their hope.
It is a paradox.
Now, as destruction on a scale never seen before descends on Gaza, he wonders if the solution -- the only solution -- is for people to just leave and save themselves.
Will they be letting Israel win?
Hossam Madhoun ( right ) and Jamal Al Rozzi ( left ) in their production of ' The Emigrants ' for their theatre company, ' Theatre for Everybody ', performed in Gaza in February 2022 ( Supplied ) This question is especially pertinent now that Madhoun himself has become a refugee.
Originally living in Gaza City with his wife Abeer, their daughter Salma and dog Buddy, Madhoun and his family left after the first Israeli expulsion orders and went to Netzarim, where Abeer ' s family lived.
Now, in Egypt, Madhoun understands more deeply the psychology of leaving his homeland, which he described two years previously in his production."It ' s like a place between earth and the sky,"he says.
They are unsure if they can continue living there because they have no legal status and there is no asylum process to speak of.
The uncertainty is a nightmare -- while they are safe for now, they never know what will come next.
His wife has just lost her job at the charity Humanity and Inclusion, which required her to live in Gaza rather than work remotely.
Without her income, they can no longer send money to family members who are still trapped in the warzone, nor can they support themselves.
This allows Israel to claim it is allowing aid to enter when, in reality, it facilitates the looting of it.
In 2014, Madhoun adapted a production of Leo Tolstoy ' s War and Peace, which advocated that nothing good can come from so much war, that human beings can not psychologically survive such an onslaught and that peace is essential for freedom."It ' s too much for one life,"Madhoun says.
Theatre for Everybody and Az Theatre are currently collaborating on a series of events run from London and Cairo called"Palestine : Repair \&amp ; Return"at P21 Gallery and online.
Middle East Eye delivers independent and unrivalled coverage and analysis of the Middle East, North Africa and beyond.
To learn more about republishing this content and the associated fees, please fill out this form.
More about MEE can be found here.

% matched lemmas: 
\end{textsample}
